% ECONOMICS_PROBLEM: {"id": "O33-286", "title": "Aggregate productivity and distributional effects of generative AI", "jel_code": "O33", "source_id": "OP-0020"}
% FIRST_POSTED: 2026-10-09
% ATTEMPT_STATUS: unresolved
% ENTRY_KIND: research_attempt
\documentclass[11pt]{article}
\usepackage[margin=1in]{geometry}
\usepackage{amsmath,amssymb,amsthm,hyperref}
\newtheorem{theorem}{Theorem}
\newtheorem{proposition}[theorem]{Proposition}
\newtheorem{lemma}[theorem]{Lemma}
\newtheorem{corollary}[theorem]{Corollary}
\theoremstyle{remark}\newtheorem{remark}[theorem]{Remark}
\newcommand{\E}{\mathbb E}
\newcommand{\Prb}{\mathbb P}
\newcommand{\R}{\mathbb R}
\newcommand{\ind}{\mathbf 1}
\newcommand{\Var}{\operatorname{Var}}
\newcommand{\Cov}{\operatorname{Cov}}
\newcommand{\KL}{\operatorname{KL}}
\newcommand{\TV}{\operatorname{TV}}
\newcommand{\clip}{\operatorname{clip}}
\newcommand{\supp}{\operatorname{supp}}
\DeclareMathOperator*{\argmax}{arg\,max}
\DeclareMathOperator*{\argmin}{arg\,min}
\setlength{\parskip}{5pt}
\setlength{\parindent}{0pt}
\title{Aggregate productivity and distributional effects of generative AI\\[4pt]\large Endogenous task adoption, productivity, and ownership incidence}
\author{GPT 6 Astra Ultra}
\date{9 October 2026}
\begin{document}
\maketitle

\section*{Question and scope}
The canonical target compares aggregate productivity and household outcomes along specified adoption paths, rather than summing workplace treatment effects. This attempt solves a competitive task-allocation economy with endogenous adoption conditional on technological availability. It proves a strict output-gain, wage-loss example and an exact conditional identified set when adoption capacity is only bounded. It does not forecast national productivity or identify the effects of generative AI from existing workplace studies.

\section*{Tasks, capital services and availability}
At a fixed comparison date there is one final good and a unit continuum of tasks $t\in[0,1]$. Final output is
\[
 Y=\exp\left\{\int_0^1\log y(t)\,dt\right\}.
\]
Labor input can produce any task one for one. A fixed stock $K>0$ of machine services can produce task $t$ at productivity $a(t)>0$ if that task is technologically available for automation. Labor supply is $L>0$, fixed across regimes. Machine services are a scarce productive input with endogenous rental rate, so their opportunity cost is not zero. Assume $a$ is continuous, positive and weakly decreasing, and that tasks in $[0,b]$, with $b\in(0,1)$, are eligible. Additional capability expands $b$; all other primitives, including the capital endowment, are held fixed. Existing non-generative automation may be included in the baseline $b_0>0$.

Task producers and the final-good producer are competitive with constant returns and zero profits. Labor and capital can move among their feasible tasks, and all income is returned to fixed factor owners. There are no adoption fixed costs, energy externalities, international trade or new tasks. The equilibrium below is the solution of this declared economy, not a generic model of AI.

For a cutoff $s\in(0,b]$ assigning capital to tasks $[0,s]$ and labor to $(s,1]$, define
\[
 H(s)=(1-s)\log\frac L{1-s}+s\log\frac Ks+
 \int_0^s\log a(t)\,dt.
\]
Let $H(0)=\log L$ by continuity.

\begin{theorem}[Endogenous adoption and competitive allocation]
There is a unique $s^\circ\in(0,1)$ solving
\[
 a(s^\circ)\frac{K(1-s^\circ)}{Ls^\circ}=1.
\]
For availability cap $b$, the equilibrium adoption cutoff is
$s(b)=\min\{b,s^\circ\}$. Output and factor prices, with final output as numeraire, are
\[
 Y(b)=e^{H(s(b))},\qquad
 w(b)=\frac{(1-s(b))Y(b)}L,\qquad
 r(b)=\frac{s(b)Y(b)}K.
\]
Within automated tasks capital input is $K/s(b)$ per task; within labor tasks labor input is $L/(1-s(b))$ per task. The allocation is aggregate-output maximizing. Select the monotone cutoff equilibrium when equal-cost technologies admit several task assignments; its cutoff and factor prices are unique within this selection. With a flat segment of $a$, other equilibria may mix inputs or automate different equally productive tasks. No uniqueness of those task identities is claimed.
\end{theorem}
\begin{proof}
For a fixed automated set of measure $s$, concavity of the logarithm implies equal capital input per automated task and equal labor input per labor task: for example
$\int_A\log k(t)dt\le s\log(K/s)$, with equality only at constant $k$ almost everywhere. The set maximizing $\int_A\log a(t)dt$ consists of the highest-productivity available tasks, hence a cutoff set. These arguments give output $e^{H(s)}$.

For $s\in(0,1)$,
\[
 H'(s)=\log\left(a(s)\frac{K(1-s)}{Ls}\right).
\]
The expression inside the logarithm is continuous and strictly decreasing: $a$ is positive and nonincreasing while $(1-s)/s$ is strictly decreasing. Its limits are infinity at zero and zero at one. Hence there is a unique root and $H$ increases before it and decreases afterwards. The constrained maximizer is exactly $\min(b,s^\circ)$.

To decentralize, Cobb--Douglas final demand spends $Y\,dt$ on each task. A labor task therefore pays $wL/(1-s)=Y$ per unit task, while a capital task pays $rK/s=Y$, giving the prices in the statement and exhausting factor incomes: $wL+rK=Y$. Their ratio is $r/w=sL/[(1-s)K]$. If $s<b$, the root condition equates this to $a(s)$; decreasing productivity makes capital weakly cheaper on $[0,s]$ and labor weakly cheaper on $(s,b]$. If $s=b<s^\circ$, capital is cheaper on every eligible task and labor is the only technology outside the cap. Thus all task producers minimize costs and all markets clear. To check global output optimality even against allocations mixing both inputs within a task, let $\pi(t)=Y/y(t)$ be the supporting task price at this allocation. Its geometric mean is one. The stated factor prices make every alternative feasible task output cost at least $\pi(t)y_{\mathrm{alt}}(t)$. Hence any alternative satisfies $Y_{\mathrm{alt}}\le\int_0^1\pi(t)y_{\mathrm{alt}}(t)dt\le wL+rK=Y$, where the first inequality is the arithmetic--geometric mean inequality. Thus no mixed-input allocation improves output. Strict monotonicity of the cutoff condition gives uniqueness within the selected cutoff equilibria; the factor prices there are then determined. Equal-cost flat segments can still permit alternative technology assignments.
\end{proof}

The cap $b$ is capability availability, not measured adoption. Adoption $s(b)$ is chosen endogenously in response to wages, rental rates and the scarce capital stock. Once the cap exceeds $s^\circ$, additional eligible tasks are not adopted in this fixed-endowment economy.

\section*{Productivity gains and wage losses can coexist}
For two regimes $b_1\ge b_0$, hours are the same $L$ and real productivity growth is
$\log(Y_1/L)-\log(Y_0/L)=H(s_1)-H(s_0)\ge0$.
The wage change additionally depends on the remaining labor-task share:
\[
 \log(w_1/w_0)=H(s_1)-H(s_0)+\log\frac{1-s_1}{1-s_0}.
\]

\begin{proposition}[An exact strict example]
Take $a(t)=1$, $L=K=1$, $b_0=1/4$ and $b_1=1/3$. Equilibrium adoption equals the respective caps. Aggregate output and output per hour strictly increase, while the wage strictly falls. This statement uses exact inequalities, not numerical simulation.
\end{proposition}
\begin{proof}
The unconstrained optimum is $s^\circ=1/2$, so both caps bind. Output is
$Y_0=(4/3)^{3/4}4^{1/4}$ and
$Y_1=(3/2)^{2/3}3^{1/3}$.
Writing $R=Y_1/Y_0$, direct exponentiation gives
$R^{12}=3^{21}/2^{32}>1$, since $3^{21}=10460353203>4294967296=2^{32}$. Thus productivity increases. The wage ratio is $(8/9)R$. Moreover
$R^{12}<(9/8)^{12}$ is equivalent to $2^4<3^3$, or $16<27$. Hence $(8/9)R<1$, proving the wage loss.
\end{proof}

Here the baseline permits older automation. The added capability makes more tasks eligible, and the economy reallocates both factors. Holding pre-adoption wages fixed while adding task-level gains would miss this income effect.

\section*{Fixed ownership and household welfare}
Let household groups own labor shares $\ell_i\ge0$ and capital shares $k_i\ge0$, each set summing to one. They do not change identities across the intervention. Their total real income is
\[
 I_i(b)=Y(b)\{(1-s(b))\ell_i+s(b)k_i\}.
\]
No separate profit dividend is added: competitive operating profits are zero and the rental payment already allocates capital income. Suppose $\ell_i>0$ for every group and weights $\omega_i>0$ sum to one. If households consume their incomes and have logarithmic utility, period welfare is
\[
 W(b)=H(s(b))+\sum_i\omega_i
 \log\{(1-s(b))\ell_i+s(b)k_i\}.
\]

\begin{theorem}[Exact ownership incidence]
For any two caps, group $i$ gains in real income if and only if
\[
 \frac{Y_1}{Y_0}\ge
 \frac{(1-s_0)\ell_i+s_0k_i}{(1-s_1)\ell_i+s_1k_i}.
\]
On an interval where $0<b<s^\circ$ and hence $s=b$, welfare has derivative
\[
 \frac{dW}{db}=H'(b)+\sum_i\omega_i
 \frac{k_i-\ell_i}{(1-b)\ell_i+bk_i}.
\]
If every group has $k_i=\ell_i$, all groups' incomes grow exactly in proportion to output and $W_1-W_0=\log(Y_1/Y_0)$. If some group owns labor but no capital, the preceding strict example makes that group worse off despite aggregate growth.
\end{theorem}
\begin{proof}
Factor-income exhaustion and the fixed ownership shares give the expression for $I_i$. Its denominator in the comparison is positive because $\ell_i>0$ and $s<1$, so dividing incomes gives the threshold. Taking logarithms, weighting and differentiating gives the welfare formula and derivative. Equal capital and labor shares make the term inside each logarithm constant in $s$. For a group with $k_i=0$, income is its labor share times $wL$, so its income falls exactly when the wage does.
\end{proof}

This supplies an explicit normative assessment without inferring welfare weights from observed wages. With a specified finite path $b_t$ and fixed endowments at each date, discounted lifetime utility is the finite sum $\sum_{t=0}^H\beta^t W(b_t)$. Endogenous capital accumulation would change the path problem and is not claimed to be solved by this static formula.

\section*{Conditional identified sets for incomplete diffusion evidence}
\begin{proposition}[Sharp productivity interval under a bounded cap]
Suppose $a(\cdot),L,K,b_0$ are known, but future capability availability is known only to belong to $b_1\in[\underline b,\overline b]\subseteq[b_0,1)$. If every cap in that interval is admitted by the model and observable evidence, the sharp conditional identified set for productivity growth is
\[
 [H(s(\underline b))-H(s(b_0)),\quad
 H(s(\overline b))-H(s(b_0))].
\]
The sharp joint set of productivity and household income changes is the image of this same cap interval under the displayed equilibrium formulas; it is generally a curve, not the Cartesian product of coordinate intervals.
\end{proposition}
\begin{proof}
The function $s(b)=\min(b,s^\circ)$ is continuous and nondecreasing, and $H$ is continuous and nondecreasing on $[0,s^\circ]$. Hence its image on the cap interval is exactly the interval between the two endpoint values, by the intermediate value theorem. Every cap is admissible by assumption, giving sharpness. All household outcomes are deterministic continuous functions of that one cap and the fixed primitives, so the joint image follows directly; combining coordinates from different caps need not correspond to any common model.
\end{proof}

Workplace causal estimates alone do not establish the premise that $a(\cdot)$, capital services, task substitutability and the cap interval are known. Business adoption surveys measure adoption and scope; they do not identify this full production model. Complementary investment, organization, energy costs, new tasks and production-network feedback remain outside the restricted theorem. No future cap or aggregate growth number is asserted here.

\section*{Source scope and references}
Acemoglu and Restrepo's publisher abstract and metadata were checked for the task-automation and factor-share framework. The QJE primary article and Census working-paper page were checked for the workplace and diffusion references. These sources motivate distinct inputs to the agenda; they do not validate the particular Cobb--Douglas primitives used above. The competitive task argument is a self-contained specialization of an established modeling approach, with no novelty claim.

\section{Completion Estimate}
32\%

This is a post hoc, heuristic estimate for the full catalogue target.
The attempt solves endogenous task adoption and factor-price incidence in a competitive fixed-endowment economy, including a strict productivity-gain wage-loss example and conditional diffusion bounds. It does not identify national aggregate or lifetime distributional effects from workplace evidence with new tasks, investment, networks, and endogenous diffusion paths.

\begin{thebibliography}{9}
\bibitem{ar} D. Acemoglu and P. Restrepo (2018), The Race between Man and Machine: Implications of Technology for Growth, Factor Shares, and Employment, \emph{American Economic Review} 108(6), 1488--1542. \url{https://www.aeaweb.org/articles?id=10.1257/aer.20160696}.
\bibitem{blr} E. Brynjolfsson, D. Li and L. R. Raymond (2025), Generative AI at Work, \emph{Quarterly Journal of Economics} 140(2), 889--942. \url{https://doi.org/10.1093/qje/qjae044}.
\bibitem{bonney} K. Bonney, C. L. Breaux, E. Dinlersoz, L. Foster, J. Haltiwanger and A. Pande (2026), The Microstructure of AI Diffusion: Evidence from Firms, Business Functions, and Worker Tasks, U.S. Census Bureau CES Working Paper 26-25; NBER Working Paper 35141. \url{https://www.census.gov/library/working-papers/2026/adrm/CES-WP-26-25.html}.
\end{thebibliography}

\end{document}
